\section{Related Work}
\label{sec:related}

\subsection{Tensor-network simulation and treewidth}

Tensor-network simulation represents amplitudes, probabilities, and expectation values as contractions of local tensors. The cost of exact contraction is controlled by the maximum intermediate tensor dimension, which can be bounded in terms of treewidth or related width parameters of a contraction graph. Markov and Shi proved that a quantum circuit whose underlying graph has treewidth \(d\) can be simulated in time polynomial in the number of gates and exponential in \(d\) \citep{MarkovShi2008}. This remains a central structural baseline: if a compositional circuit is genuinely tree-like or has logarithmic treewidth, it is classically tractable by standard contraction methods.

The present paper does not improve the general Markov--Shi theorem for arbitrary dense tensors. Instead, it refines the simulation question for tensor networks whose tensors have an additional stabilizer/resource structure. The relevant comparison is between the ordinary dense separator width \(\wTN\) and an effective stabilizer-compressed width \(\weff\). If these widths coincide, the separator-local theorem is not automatically better than ordinary contraction.

Dense treewidth and stabilizer tableau simulability can disagree sharply. A two-dimensional graph-state tensor network has a dense contraction graph with grid treewidth, but the same state is a stabilizer object with a polynomial-size tableau description. The separation examples in this paper exploit precisely this mismatch: dense tensor-network bounds see a high-width graph, while SLMS first certifies the free stabilizer region and charges only for dense active boundary variables and non-marginalized resource choices. This is not a lower bound against stabilizer-aware classical algorithms; it is a separation between dense tensor-network width parameters and stabilizer-compressed separator-local parameters.

Treewidth is also not the only classical tensor-network control parameter. Matrix-product-state and tree-tensor-network simulations can be efficient when the relevant bond dimensions, Schmidt ranks, or truncation errors remain controlled \citep{Vidal2003,Schollwoeck2011}. Modern tensor-network simulation practice further exploits geometry, approximate contraction, belief propagation, symmetries, and problem-specific correlation structure. For example, Tindall, Fishman, Stoudenmire, and Sels simulated IBM's Eagle kicked-Ising experiment using a tensor-network approach adapted to the heavy-hex geometry and approximately contracted with belief propagation \citep{TindallEtAl2024}. The present paper is therefore not positioned as a general replacement for optimized tensor-network algorithms. Its exact dense-width comparison is against dense treewidth-style contraction bounds, while its positive result exploits an orthogonal stabilizer/quasiprobability certificate.

\subsection{Relation to junction-tree inference and focused-width simulators}

The active-child recurrence is best viewed as a junction-tree-style collect recurrence lifted to stabilizer-compressed quasiprobability tensor networks. Classical junction-tree and local-computation algorithms combine local factors and marginalize variables along a tree decomposition or hypertree \citep{LauritzenSpiegelhalter1988,ShaferShenoy1990}; bucket elimination gives an equivalent elimination-order perspective \citep{Dechter1999}. SLMS inherits this message-passing skeleton. It also inherits the signed-sampling logic of quasiprobability simulation \citep{PashayanWallmanBartlett2015,PashayanBartlettGross2020}: local non-free tensors are expanded into finite free dictionaries, sampled from absolute weights, and reweighted by signs. Finally, it inherits stabilizer compression from Gottesman--Knill/tableau simulation.

The contribution is not the recurrence alone. The contribution is a checkable sufficient certificate combining these ingredients for locally marginalizable structured tensor networks. The certificate records finite local free dictionaries, bag costs \(\xi_b\), local coefficient expansion factors \(\kappa_b\), stabilizer-compressed boundary width \(\weff\), and active-child declarations \(\Act(b)\). Local marginalizability gives operational semantics to the statement that a resource label has been consumed before it crosses a separator. This is information that a global magic burden does not express.

Focused tree-width, focused rank-width, ZX partitioning, and stabilizer tensor-network methods are direct competitors rather than special cases of SLMS. SLMS remains useful if focused-width methods exist because it reports a different, local-marginalization certificate: it can explain why a proposed compositional ansatz is easy, identify the bags where the certificate fails, and provide a scalar additive estimator when the certificate succeeds. The paper does not claim that \(\weff+2\Lammsg\) dominates focused tree-width, focused rank-width, low-rank-width ZX parameters, or stabilizer tensor-network costs.

\subsection{Focused width, ZX partitioning, and stabilizer tensor networks}

Recent work substantially sharpens the classical simulation landscape at the intersection of tensor networks, stabilizer decompositions, and ZX-calculus. Codsi and Laakkonen introduce focused tree-width and focused rank-width, together with simulation algorithms that unify graph measures and stabilizer decompositions \citep{CodsiLaakkonen2026}. Kuyanov and Kissinger give a ZX-calculus simulation method for low-rank-width graph-like ZX diagrams with complexity governed by rank decompositions \citep{KuyanovKissinger2026}. These papers are especially relevant because they also combine graph structure and non-Clifford resource, and their parameters may dominate \(\weff+2\Lammsg\) on many instances.

ZX-calculus cutting and decomposition methods provide another strong comparator. Sutcliffe's \(k\)-partitioning method and Sutcliffe--Kissinger's ZX cutting and parametric rewriting methods use diagram rewriting, partitioning, and shared reductions across related branches to improve strong simulation of Clifford+\(T\) and related circuits \citep{Sutcliffe2024KPartition,SutcliffeKissinger2024Cutting,SutcliffeKissinger2025Parametric}. Stabilizer tensor-network approaches also combine tableau structure with tensor-network ideas: Masot-Llima and Garc\'ia-S\'aez introduce stabilizer tensor networks as a universal simulator on a stabilizer-state basis \citep{MasotLlimaGarciaSaez2024}, and Nakhl and collaborators extend this direction using magic-state injection \citep{NakhlEtAl2025}.

The present paper does not claim to beat these methods. Its contribution is narrower: it gives a certificate-based message recurrence for locally marginalizable compositional tensor networks and records the resulting parameter \(\Lammsg\). A complete empirical comparison must compute focused tree-width or focused rank-width, run relevant ZX-inspired partitioning baselines where available, and compare against stabilizer tensor-network implementations. Without such evidence, SLMS is best interpreted as a diagnostic and restricted theorem rather than a state-of-the-art simulator.

\subsection{Stabilizer simulation and magic resources}

The Gottesman--Knill theorem shows that stabilizer circuits can be simulated efficiently; tableau simulation gives an explicit polynomial-time data structure for this closure \citep{AaronsonGottesman2004}. Modern simulation methods for near-stabilizer circuits extend this by expanding non-Clifford states, gates, or channels into stabilizer components. Bravyi and Gosset gave improved simulation algorithms for Clifford+\(T\) circuits whose cost scales exponentially in a function of the number of \(T\) gates \citep{BravyiGosset2016}. Bravyi, Browne, Calpin, Campbell, Gosset, and Howard developed a systematic theory of stabilizer rank and approximate stabilizer rank and applied it to low-rank stabilizer decompositions \citep{BravyiEtAl2019}.

Quasiprobability methods give a different operational picture. Pashayan, Wallman, and Bartlett estimate outcome probabilities by sampling from signed or quasiprobability representations, with sample complexity governed by an \(\ell_1\)-type negativity \citep{PashayanWallmanBartlett2015}. Howard and Campbell related robustness of magic to classical simulation overhead in a resource-theoretic setting \citep{HowardCampbell2017}. Seddon, Regula, Pashayan, Ouyang, and Campbell sharpened these links through tighter magic monotones and improved simulation algorithms \citep{SeddonEtAl2021}.

Our separator-local burden uses these operational quantities rather than replacing them. The primitive local costs \(\xi_b\) are certified free-tensor quasiprobability \(\ell_1\) costs; stabilizer extent, robustness, and dyadic negativity are relevant when they instantiate or upper-bound that cost. The new question is how these local costs compose over a tree decomposition. In the present formulation this composition is made explicit through a message recurrence for locally marginalizable networks, rather than through an existential support-set definition.

\subsection{Why this is not just Gottesman--Knill plus local sampling}

A purely stabilizer network is already classically easy by Gottesman--Knill simulation. A global quasiprobability simulator already handles non-Clifford resources by expanding them into free components, but it pays for the global product of local \(\ell_1\) costs. The object studied here is the intermediate case in which resource tensors are present but their quasiprobability labels may be locally consumed before they load every separator of a compositional graph.

The new ingredient is therefore not stabilizer simulation alone, nor local sampling alone, but a certified tree-decomposition recurrence. The recurrence combines stabilizer compression, local quasiprobability costs, active-child certificates, and local marginalization maps to prove that inactive resource randomness does not multiply into ancestor messages. This is a restricted statement, but it yields a computable parameter that can separate global resource burden from message burden and dense treewidth from stabilizer-compressed effective width.

\subsection{Stabilizer R\'enyi entropy and diagnostics}

Stabilizer R\'enyi entropies and related nonstabilizerness diagnostics provide computable probes of magic. Leone, Oliviero, and Hamma introduced stabilizer R\'enyi entropy for pure states and studied its resource-theoretic properties and many-body applications \citep{LeoneOlivieroHamma2022}. Haug and Piroli analyzed stabilizer entropies and nonstabilizerness monotones, including limitations of monotonicity in certain regimes \citep{HaugPiroli2023}. Leone and Bittel later proved monotonicity results for stabilizer entropies in magic-state resource theory under specified conditions \citep{LeoneBittel2024}.

In this paper, \(M_2\), linear stabilizer entropy, and related quantities are treated as diagnostics for benchmark analysis. They are not used as the primary operational simulation overhead unless a rigorous conversion to stabilizer extent, robustness, or quasiprobability norm is available for the tensors under consideration.

\subsection{QNLP and compositional text circuits}

Categorical compositional distributional semantics was introduced by Coecke, Sadrzadeh, and Clark as a mathematical framework combining grammatical composition with distributional meaning \citep{CoeckeSadrzadehClark2010}. Zeng and Coecke proposed quantum algorithms for compositional natural language processing \citep{ZengCoecke2016}. Meichanetzidis and collaborators developed a near-term QNLP pipeline mapping grammatical diagrams to circuits \citep{MeichanetzidisEtAl2020}, and early practical QNLP experiments on NISQ hardware were reported by Lorenz and collaborators \citep{LorenzEtAl2021}. lambeq provides a high-level Python library for converting sentences to string diagrams, tensor networks, and quantum circuits \citep{KartsaklisEtAl2021}.

QDisCoCirc extends this compositional viewpoint to discourse-level text processing. Laakkonen, Meichanetzidis, and Coecke analyze quantum algorithms for compositional text processing and establish conditional hardness for certain question-answering tasks under explicit assumptions \citep{LaakkonenMeichanetzidisCoecke2024}. Related tensor-network models for sequence processing emphasize structured, often tree-like, connectivity and quantum-inspired contraction \citep{HarveyYeungMeichanetzidis2025}. These works motivate the family of graphs studied here. The present paper contributes a simulation and resource-localization framework, not a new QNLP model.

\subsection{Dequantization and QML simulability}

The dequantization literature shows that proposed quantum speedups may disappear when classical algorithms exploit the same input assumptions, low-rank structure, or tensor-network geometry. Tang's quantum-inspired recommendation algorithm \citep{Tang2019}, subsequent low-rank dequantization frameworks \citep{ChiaEtAl2020}, and broader discussions of dequantization in quantum machine learning \citep{Tang2022} are directly relevant as methodological caution. Recent tensor-network analyses of variational quantum machine-learning models further show that expressive quantum models may define classically tractable function classes under structural restrictions \citep{ShinTeoJeong2024}.

For this reason, this paper is framed as a classical simulation theorem plus a resource-localization framework. It is not an advantage claim. If a family of compositional quantum language circuits has small width and small separator-local magic, the correct conclusion is that the family is classically simulable by the present methods or by nearby tensor-network/stabilizer methods.
